% !PRINTED 648
\weberpage{648}
%<<CONT
are conjugate imaginary; for, first, their $n$th powers are conjugate
imaginary, and hence they themselves are so, at first up to an $n$th root of unity as
factor. But since their product is real, this $n$th root of unity
must be $= 1$. Since, moreover,
\[
r_\nu \equiv - r_{\nu + \frac{n-1}{2}} \pmod{n}
\]
formula (11) shows that in this case the roots $\xi_z$ are all real.

With this, then, the theorem is proved for a real field $\mathfrak{K}$:

\gesperrt{If the cyclic equation whose roots are the $k_0$, $k_1 \ldots k_{n-2}$ has real
roots, then of the roots $\xi$ one is real and the rest imaginary; if this
cyclic equation has imaginary roots, then all the $\xi$ are real.}

\begin{center}
§. 188.

\gesperrt{Metacyclic Equations of the Fifth Degree.}
\end{center}

Through the theorems of the preceding section the roots of a metacyclic
equation of prime degree have been brought into a generally valid form, which permits
all these quantities to be actually formed, if one presupposes a knowledge of the roots of cyclic
equations. Whole systems of cyclic equations of every degree are furnished to us,
e.g., by the theory of the division of the circle, whose roots are the cyclotomic periods. Thus
we can, for example, form arbitrarily many
metacyclic equations for the field of rational numbers. That all these equations in the field of
rational numbers are contained among them is a very remarkable theorem due to \gesperrt{Kronecker},
which we shall come to know in the second volume. Here, as an
application and illustration of what precedes, we wish further to treat the special problem
of finding, in an arbitrary field $\mathfrak{K}$, the roots of all metacyclic
equations of the fifth degree.

For this it is necessary and sufficient that we determine generally the roots $k_0, k_1, k_2, k_3$ of a
cyclic equation of the fourth degree, and indeed under the
assumption that \gesperrt{$k_0, k_1, k_2, k_3$ are different from one another and none
of them is zero}.
